\documentclass[a4paper,12pt]{article}

% Set margins
\usepackage[hmargin=2.5cm, vmargin=2.5cm]{geometry}

\frenchspacing

% Language packages
\usepackage[utf8]{inputenc}
\usepackage[T1]{fontenc}
\usepackage[magyar]{babel}

% AMS
\usepackage{amssymb,amsmath}

% Graphic packages
\usepackage{graphicx}

% Colors
\usepackage{color}
\usepackage[usenames,dvipsnames]{xcolor}

% Enumeration
\usepackage{enumitem}

% Multiple rows for the journal title heading
\usepackage{multirow}

% Links
\usepackage{url}
\usepackage{hyperref}

\usepackage{fancyhdr}

\fancyhead[LO, LE]{Pi Matematikai, Fizikai és Informatikai Folyóirat}
\fancyhead[RO, RE]{1. kötet, Matematikai füzet, 2024.}

\usepackage{amsthm}

%\newtheorem{definition}{Definíció}[section]
%\newtheorem{notation}{Jelölés}[section]
\newtheorem{theorem}{Tétel}[section]
\newtheorem{lemma}[theorem]{Lemma}
\newtheorem{proposition}[theorem]{Propozíció}%[section]

\theoremstyle{definition}
\newtheorem{definition}[theorem]{Definíció}
\newtheorem{notation}[theorem]{Jelölés}
\newtheorem{problem}[theorem]{Probléma}%[section]
\newtheorem{corollary}[theorem]{Következmény}%[section]
\newtheorem{remark}[theorem]{Megjegyzés}%[section]
\newtheorem{example}[theorem]{Példa}%[section]
\newtheorem{exercise}[theorem]{Feladat}%[section]
%\newtheorem{solution}[theorem]{Megoldás}%[section]
\newtheorem{conjecture}[theorem]{Sejtés}%[section]

\renewenvironment{proof}[1][]{\noindent\textit{Bizonyítás} }{$\qed$}
\newenvironment{solution}[1][]{\noindent\textit{Megoldás} }{$\qed$}

\begin{document}

\pagestyle{fancy}

\title{\sc A cikk címe}

% Kérem írja be a szerzők nevét, hovatartozását, e-mail címét a megfelelő helyre. 
\author{
	\textbf{Első Szerző} \\
	{\normalsize Affiliáció 1} \\
	{\normalsize \href{mailto:első.szerző@email.com}{első.szerző@email.com}}
\and
	\textbf{Második Szerző} \\
	{\normalsize Affiliáció 2} \\
	{\normalsize \href{mailto:második.szerző@email.com}{második.szerző@email.com}}
\and
	\textbf{Harmadik Szerző} \\
	{\normalsize Affiliáció 3} \\
	{\normalsize \href{mailto:harmadik.szerző@email.com}{harmadik.szerző@email.com}}
}

%Az utolsó verzió dátuma
\date{\textit{2024. szeptember 24.}}

\noindent \begin{tabular}{lp{13cm}}
	\multirow{3}{*}{\includegraphics[width=22mm, height=22mm]{jmathpi_logo.pdf}}
	& Pi Matematikai, Fizikai és Informatikai Folyóirat \\
	& Miskolci Egyetem \\
	& \\
	& 1. kötet, Mat. füzet, 1-3. oldal, 2024. \hfill \href{https://doi.org/10.1000/jmathpi-1}{DOI: 10.1000/jmathpi-????} \\
\end{tabular}

\bigskip

\hrule

{\let\newpage\relax\maketitle}

\begin{abstract}
Absztrakt szövege.
\end{abstract}

{\def\thefootnote{}\footnotetext{Kulcsszavak: mat., AMS Subject Classification: AB123, CD567}}

\section{Bevezetés}

A bevezetés szövege. 

\section{A második fejezet címe}

Kérjük, az alábbi környezeteket használja:

\begin{definition}
A definíció szövege
\end{definition}

\begin{notation}
A jelölés szövege.
\end{notation}

\begin{theorem}\label{theorem:Key1}
A tétel szövege. Kérjük, hogy így használja a szövegben az olyan matematikai formulákat, mint például az $\frac{a}{b}$. 
\par 
\[
\text{Számozás nélküli matematikai egyenlet szövege}
\]
\begin{equation}\label{equation:Key2}
\text{Számozott matematikai egyenlet szövege, mint pl.} \dfrac{a}{b}
\end{equation}
Kérjük használja a matematikai formulákat olyan matematikai környezetben, mint például a \ref{equation:Key2} egyenletben. 
\end{theorem}

\par A \ref{theorem:Key1} Tétel alapján...

\begin{lemma}
A lemma szövege.
\end{lemma}

\begin{proposition}
A propozíció szövege. 
\end{proposition}

\begin{problem}
A probléma szövege. 
\end{problem}

\begin{corollary}
A következmény szövege
\end{corollary}

\begin{remark}
A megjegyzés szövege. 
\end{remark}

\begin{example}
A példa szövege. 
\end{example}

\begin{exercise}
A feladat szövege. 
\end{exercise}

\begin{solution}
A megoldás szövege. 
\end{solution}

\begin{conjecture}
A sejtés szövege. 
\end{conjecture}

\begin{proof}
A bizonyítás szövege. 
\end{proof}

\section{A harmadik fejezet címe}

\begin{table}[!ht] % vagy [H] az elúszás kikapcsolásakor
  \centering
  \caption{A táblázat felirata.}\label{table:Key3}
  \medskip
  \begin{tabular}{|c|c|}
    \hline
    1 & 2 \\
    \hline
    3 & 4 \\
    \hline
  \end{tabular}
\end{table}

Az \ref{table:Key3}. táblázat alapján\dots

\begin{figure}[!ht] % vagy [H] az elúszás kikapcsolásakor
  \centering
  \includegraphics[width=3cm]{ME_Logo.png}
  \caption{A kép felirata.}\label{figure:Key4}
\end{figure}

Az \ref{figure:Key4}. ábra alapján\dots

Kérjük használjon különböző listákat, például
\begin{itemize}
\item rendezetlen lista,
\item rendezett lista,
\item hierarchikus lista.
\end{itemize}

\section{A negyedik fejezet címe}%stb.

Hivatkozzon könyvekre \cite{erdos2013}, illetve cikkekre \cite{aczel1983}.

\paragraph{Köszönetnyilvánítás.} A köszönetnyilvánítás szövege.

\bibliographystyle{plaindin}
\bibliography{references}

\end{document}
